% Pursuing Stacks (cotes 134-2 to 134-8): archive cote 134-6, site page 38 (PDF page 39), typescript page 356.
% typescript. Draft compared: ps5.tex:3207-3289.
% AI-transcribed by Claude Opus 5.5 (claude-opus-5-5, adaptive thinking, effort high), Message Batches API, batch ps3 (msgbatch_01XdFmuVVBAcgc7EFxdtCqsj), 2026-11-12; reviewed against both supplied scan halves; see CHECK_DECISIONS.tsv.
\scanpage{134-6/38}
\mspage{356}

a strict test category and $A^\wedge$ is $\underline{W}$-modelizing,  it follows that $A^\wedge$ is
even a \uline{canonical modelizer} (with respect to $\underline{W}$ ), \tikz[baseline=-.2ex]{\draw (0,.6ex)--(.35ex,.3ex);} i.e.  defined in
terms of the $\underline{W}$-asphericity structure associated to the ``canonical''
homotopy structure $h'' = h_{A^\wedge}$ on $A^\wedge$ (cf. prop.2  (ii)  p. 345).  This
for the time being,  together with the modelizers (Cat) and (Spaces),
are the main examples we got of canonical modelizers.  Presumably
stacks should give another sizeable bunch of canonical modelizers,
not of the type $A^\wedge$.

\begin{center}\shortstack{*\qquad *\\{}*}\end{center}

We still have to say a word about morphisms \inserted[typed]{between} \xout{of} contractors $A,B$ .
T\over[typed]{h}{h}e first thing that comes to my mind is that this should be a functor
\hand{(0)} \[ f\colon A \hand{\longrightarrow} B \]
such that the corresponding functor
\hand{(1)} \[ f^*\colon B^\wedge \hand{\longrightarrow} A^\wedge \]
should be compatible with the homotopy structures,  which can be
expressed\hand{,} as we know\hand{,} in manifold ways,  the most natural one here
being the following two
(2) \[ f^*(\over{A}{B}) \hand{\subset} \over{B}{A}^\wedge_c \]
or
(3) \[ f^*(\over{A}{B}^\wedge_c) \hand{\subset} \over{B}{A}^\wedge_c \quad , \]
which are both implie\over[typed]{e}{d} by the apparently weaker one
(4) \[ f^*(\over{A}{B}) \hand{\subset} \over{B}{A}^\wedge_{\underline{W}_0} \overset{\text{def}}{=} \text{ set of 0-connected objects of } A^\wedge, \]
and equivalently still,  \over{t}{a}\over{\unsure{o}}{s} $\over{B}{A}^\wedge_c \hand{\subset} \over{B}{A}^\wedge_{\underline{W}} \hand{\subset} A^\wedge_{\underline{W}_0}$   (where $\underline{W}$  is a basic loca%
% NOTE: the right edge of the scan cuts the line after "loca"; a hyphen there would be hidden
$\underset{\xout{\illegible{3}}}{\text{lis}}$er satisfying L\tikz[baseline=(lo.base)]{\node[inner sep=0pt] (lo) {o};\draw ([xshift=-.4ex,yshift=.7ex]lo.north) .. controls ([yshift=-.1ex]lo.north) .. ([xshift=.4ex,yshift=.7ex]lo.north);\draw ([xshift=-.15ex,yshift=-.55ex]lo.south)--([xshift=.12ex,yshift=-.45ex]lo.south);}c 4)),  to the condition
(5) \[ f^*(B) \hand{\subset} A^\wedge_{\underline{W}} \quad (\overset{\text{def}}{=} \text{ \over{\illegible{1}}{s}et of } \underline{W}\text{-aspheric objects of } A^\wedge) \ . \]
Thus the condition for $f$ to be a ``morphism of contractors'' just boils
down to \inserted[typed]{the long familiar} $\underline{W}$-\uline{asphericity} of $f$,  and implies the following \inserted[typed]{relation} \xout{condition},  ap-%
parently stronger than (5):
(5') \[ B^\wedge_{\underline{W}} = (f^*)^{-1}(A^\wedge_{\underline{W}}) \ . \]
This comes almost as a surprise (after \inserted[typed]{a} four day's  interruptio\over[typed]{g}{n} in con-%
tact with the stuff\hand{!}) - but it occurs \inserted[typed]{to me now} that we got already a more general
statement with prop.4 of section 95 (p.350),  which includes the situa-%

\margin{38}

% Visual collation: both supplied image halves inspected; see CHECK_DECISIONS_134_7_8.tsv.
